\documentclass[12pt]{article} 

\usepackage[top=1.5cm, bottom=1.5cm, left=1.5cm, right=1.5cm]{geometry}
\usepackage{color,graphicx,amsmath,mathtools,amssymb,amsthm}
\pagestyle{plain}

\begin{document}  
\pagestyle{empty}
 




%\vskip0.5in
%
%
%\noindent {\small This document is property of Northeastern University.  Unauthorized distribution of the content is not~permitted.}
%
%\vskip0.5cm







Let $c$ be the number of chocolate chip cookies and $m$ be the number of caramel cookies. The following must be true for any selection of cookies:
$c+m=8$, $c\geq 2$, and $m\geq 2$.\\
\begin{enumerate}
\item
If the cookies of the same flavor are indistinguishable, then the only selection that matters is the number of chocolate chip cookies and caramel cookies 
that are selected. Thus, each pair of cookies $(c,m)$ pertains to one such selection. Which are the following:\\
$$(2,6), (3,5), (4,4), (5,3), (6,2)$$\\
Thus the formula for the total number of selections is:\\
$$\sum\limits_{c=2}^{6} 1 $$\\
And the total number of possible selections is $5$.\\
\item
If every cookie is distinguishable and $c$ amount of chocolate chip cookies are chosen, then there $8-c$ caramel cookies must be chosen. In other words, 
there are 10 chocolate chip cookies to choose from and $c$ choices to make. Moreover, there are 12 caramel cookies to choose from and $8-c$ choices to make. 
Thus, the formula for the total number of selections is:\\
$$\sum\limits_{c=2}^{6} \binom{10}{c} \cdot \binom{12}{8-c}$$\\
Which evaluates to $309870$, which is the total numbe of possible selections.\\



\end{enumerate} 




\end{document}